\documentclass[10pt,a4paper]{amsart}
\usepackage[margin=19mm]{geometry}
\usepackage{amsmath,amssymb,mathtools}
\usepackage[scr=rsfs,cal=euler]{mathalfa}
\usepackage{xfrac}
\usepackage[unicode,hidelinks,bookmarks=false]{hyperref}
\newcommand{\ie}{i.e.\ }
\newcommand{\cf}{cf.\ }
\newcommand{\resp}{respectively\ }
\newcommand{\Subsection}{\subsection}
\providecommand{\nobreakdash}{\nobreak\hbox{-}}
\setlength{\emergencystretch}{2em}
\multlinegap0pt
\usepackage{bm}
\usepackage{bbm}
\usepackage{dsfont}
\usepackage{xspace}
\usepackage{cancel}
\usepackage{mathtools}
\usepackage{amsthm}
\makeatletter
\@ifundefined{theorem}{\newtheorem{theorem}{Theorem}}{}
\@ifundefined{lemma}{\newtheorem{lemma}[theorem]{Lemma}}{}
\makeatother
\title[Non-uniform pairwise cross $t$-intersecting families]{Non-uniform pairwise cross $t$-intersecting families}
\author{Yongjiang Wu, Yongtao Li, Tingzeng Wu, Lihua Feng}
\date{Source: arXiv:2508.19725v3 (CC BY 4.0)}
\begin{document}
\maketitle

\noindent\emph{Source and licence.} Non-uniform pairwise cross $t$-intersecting families. Version arXiv:2508.19725v3 (CC BY 4.0), distributed under the Creative Commons Attribution 4.0 International licence, as recorded in the arXiv licence metadata for this exact version and shown on the abstract page https://arxiv.org/abs/2508.19725v3. Full rights notice: licenses/arxiv-2508.19725-CC-BY-4.0.txt.\par\smallskip

\noindent\textit{Editorial scope.} Counted unit: the Lemma whose source label is \texttt{le22} in the article (source0.tex lines 403--416, source label \texttt{le22}), delivered together with the article's own complete original proof of it (source0.tex lines 419--566, one \texttt{proof} environment of 148 lines). No mathematics is written, repaired or re-derived: statement and proof are the article's own text line for line. Required context retained uncounted: le1 (lines 311--316); le2 (lines 386--389). Every definition, display and label of those lines is retained unchanged. Omitted from this package and not redistributed: 1--310, 317--385, 390--402, 417--418, 567--1361. Nothing inside a retained excerpt is omitted or altered: every retained body is delivered exactly as the article's lines are written, so no omission receipt of any kind accompanies this package. Citations: the retained excerpts carry no citation command at all, so no bibliography entry of the article is transcribed and no reference is redistributed. Reference adaptation: none was needed and none was made; every reference inside the delivered unit resolves inside it, so no unresolved reference or citation is produced. Printed slips kept literal: no printed word, symbol, hypothesis or number of the counted statement or of its proof was altered. Layout and class: the article's own class is \texttt{article} with packages including amsmath, amsfonts, amsthm, amssymb, mathtools, mathrsfs, graphicx, color, latexsym, tikz, calc, tikz-cd, epstopdf, enumerate; the wrapper is the lane's pinned \texttt{amsart} editorial wrapper at 10pt on A4, which restates the theorem-like environments the retained text needs and adds the article's own macro definitions that the retained text needs (none), with extra wrapper packages (bm, bbm, dsfont, xspace, cancel, mathtools). The delivered numbering is the unit's own. Assets: no retained span contains a figure environment, a graphics-inclusion command, a PDF-page inclusion, a table or a TikZ picture; no image file is copied into this package, the package declares no asset dependency and no image file is redistributed with it. Licence: version arXiv:2508.19725v3 of this article is distributed under the Creative Commons Attribution 4.0 International licence, as recorded in the arXiv licence metadata for this exact version and shown on the abstract page https://arxiv.org/abs/2508.19725v3. A full rights notice is delivered at licenses/arxiv-2508.19725-CC-BY-4.0.txt. Characters: every retained excerpt is pure ASCII, so the delivered engine, XeLaTeX with its default Latin Modern text font, reports no missing character.
\begin{lemma}\label{le1}
Let $\mathcal{A}\subseteq 2^{[n]}$ be a monotone shifted family with its extent $\ell\geq  2$, generating family $\mathcal{G}$ and boundary generating family $\mathcal{G}[\ell]$.  For any families $\mathcal{B}, \mathcal{C}\subseteq \mathcal{G}[\ell]$, we denote $\mathcal{D}=(\mathcal{G}\setminus(\mathcal{B}\cup\mathcal{C}))\cup\{C\setminus\{\ell\}: C\in\mathcal{C}\}$. Then
$$
{ |\langle\mathcal{D}\rangle|=|\mathcal{A}|-|\mathcal{B}\backslash\mathcal{C}|2^{n-\ell}+|\mathcal{C}|2^{n-\ell}.}
$$
\end{lemma}

\begin{lemma}\label{le2}
 Let $\vec{\mathcal{A}} = (\mathcal{A}_1, \ldots, \mathcal{A}_m) \in (2^{[n]})^{m}$ be  a monotone shifted pairwise cross $t$-intersecting sequence with  generating sequence  $\vec{\mathcal{G}} = (\mathcal{G}_1, \ldots, \mathcal{G}_m)$. Let $\ell$ be the extent of $\vec{\mathcal{A}}$. If 
$A\in \mathcal{G}_i[\ell]$ and  $B\in \mathcal{G}_j[\ell]$  satisfy $|A\cap B|=t$ for some $i\neq j\in [m]$, then $A\cup B=[\ell]$ and $|A|+|B|=t+\ell$.
\end{lemma}

\begin{lemma}\label{le22}
  Let $\vec{\mathcal{A}} = (\mathcal{A}_1, \ldots, \mathcal{A}_m) \in (2^{[n]})^{m}$ be  a non-empty monotone shifted pairwise cross $t$-intersecting sequence with  generating sequence  $\vec{\mathcal{G}} = (\mathcal{G}_1, \ldots, \mathcal{G}_m)$. Let $\ell$ be the extent of $\vec{\mathcal{A}}$ with $\ell> t$.
Assume that $\vec{\mathcal{G}}[\ell]^{(t)}=(\emptyset,\ldots,\emptyset)$. Then 
the following statements hold.
\begin{enumerate}
    \item[(1)] If there exist some $u\neq v\in [t,\ell]$ with  $u+v=\ell+t$ such that exactly one of  $\vec{\mathcal{G}}[\ell]^{(u)}$ and $\vec{\mathcal{G}}[\ell]^{(v)}$ is not equal to  $(\emptyset,\ldots, \emptyset)$, then there exists a non-empty monotone pairwise cross $t$-intersecting sequence $\vec{\mathcal{B}}$  with its extent at most $\ell$ such that $\|\vec{\mathcal{B}}\|>\|\vec{\mathcal{A}}\|$.
    
    \item[(2)]  If  $2\nmid \ell+t$, then there exists a non-empty monotone pairwise cross $t$-intersecting sequence $\vec{\mathcal{B}}$  with its extent less than $\ell$  such that $\|\vec{\mathcal{B}}\|\geq \|\vec{\mathcal{A}}\|$.

    
    \item[(3)]   If  $2\mid \ell+t$ and for all pairs $u\neq v\in [t,\ell]$ with  $u+v=\ell+t$,  either $\vec{\mathcal{G}}[\ell]^{(u)}\neq (\emptyset,\ldots, \emptyset)$ and $\vec{\mathcal{G}}[\ell]^{(v)}\neq (\emptyset,\ldots, \emptyset)$, or $\vec{\mathcal{G}}[\ell]^{(u)}=\vec{\mathcal{G}}[\ell]^{(v)}=(\emptyset,\ldots, \emptyset)$, then there exists a non-empty monotone shifted pairwise cross $t$-intersecting sequence $\vec{\mathcal{B}}$  with its extent at most $\ell$  such that $\|\vec{\mathcal{B}}\|\geq \|\vec{\mathcal{A}}\|$. Moreover, let $\vec{\mathcal{G}}^*$ be the generating sequence of $\vec{\mathcal{B}}$.
If the extent of  $\vec{\mathcal{B}}$ is $\ell$, then  $\vec{\mathcal{G}}^*[\ell]= \vec{\mathcal{G}}[\ell]^{(\frac{\ell+t}{2})}$.
\end{enumerate}
\end{lemma}

\begin{proof}
(1)  By symmetry, we may assume that $\vec{\mathcal{G}}[\ell]^{(u)}\neq (\emptyset,\ldots, \emptyset)$ and $\vec{\mathcal{G}}[\ell]^{(v)}= (\emptyset,\ldots, \emptyset)$.
In view of $\vec{\mathcal{G}}[\ell]^{(t)}=(\emptyset,\ldots,\emptyset)$, we have $u>t$.
Recall that for a family $\mathcal{A}\subseteq 2^{[n]}$ and an integer $\ell \in [n]$, we denote $\mathcal{A}(\ell) = \{A\backslash \{\ell\}: \ell\in A\in \mathcal{A}\}$.
Let  
$$ \vec{\mathcal{G}}'= \vec{\mathcal{G}}- \vec{\mathcal{G}}[\ell]^{(u)}+\vec{\mathcal{G}}[\ell]^{(u)}(\ell),$$
where we write $\vec{\mathcal{G}}[\ell]^{(u)}(\ell)= \left(\mathcal{G}_1[\ell]^{(u)}(\ell),\ldots, \mathcal{G}_m[\ell]^{(u)}(\ell) \right)$. 
By Lemma \ref{le2}, we know that $ \langle\vec{\mathcal{G}}'\rangle$ is a non-empty
pairwise cross $t$-intersecting sequence. In addition, we have $\|\langle\vec{\mathcal{G}}'\rangle\|>\|\langle\vec{\mathcal{G}}\rangle\|=\|\vec{\mathcal{A}}\|$. Setting $\vec{\mathcal{B}}=\langle\vec{\mathcal{G}}'\rangle$ yields the desired result.




(2)   
Since $\vec{\mathcal{G}}[\ell]\neq (\emptyset,\ldots, \emptyset)$  and $2\nmid \ell+t$, it follows  that there exist $u\neq v\in [t,\ell]$ with  $u+v=\ell+t$ such that at least one of $\vec{\mathcal{G}}[\ell]^{(u)}$ and $\vec{\mathcal{G}}[\ell]^{(v)}$ is not equal to  $(\emptyset,\ldots, \emptyset)$. 
For every such pair,  we choose exactly one of the two boundary layers having
the larger norm; in case of equality, choose either one.
Let $\mathcal C$ be the union of all chosen boundary layers.
Then $\|\vec{\mathcal{G}}[\ell]- \mathcal C\|\leq \|\mathcal C\|$ and $ \|\mathcal C\|>0$.
Define 
$$
\vec{\mathcal{G}}^*
=
\vec{\mathcal{G}}-\vec{\mathcal{G}}[\ell]+
\{C\setminus\{\ell\}:C\in\mathcal C\},\qquad \vec{\mathcal{B}}=\langle\vec{\mathcal{G}}^*\rangle.
$$
We also note that every component $\mathcal G_k^{*}$ is non-empty.
Indeed, this is immediate if
$\mathcal G_k\setminus\mathcal G_k[\ell]\ne\emptyset$.
Otherwise every generator in $\mathcal G_k$ contains $\ell$.
Shiftedness then forces $\mathcal G_k=\{[\ell]\}$. Since
$\vec{\mathcal G}[\ell]^{(t)}=(\emptyset,\ldots,\emptyset)$, the
$\ell$-layer is selected from the complementary pair $\{t,\ell\}$.
Consequently, $[\ell-1]\in\mathcal G_k^{*}$.
By  Lemma \ref{le2}, $\vec{\mathcal{B}}$ is pairwise cross $t$-intersecting. 
Morover, Lemma \ref{le1} implies
\begin{align*}
\vec{\mathcal{B}}=\|\langle \vec{\mathcal{G}}^*\rangle\| =\|\vec{\mathcal{A}}\|-\|\vec{\mathcal{G}}[\ell]-\mathcal C\|2^{n-\ell}+\|\mathcal C\|2^{n-\ell}\geq \|\vec{\mathcal{A}}\|.
\end{align*}
Thus, we obtain  a  non-empty pairwise cross $t$-intersecting sequence $\vec{\mathcal{B}}$  with its extent less than $\ell$ and $\|\vec{\mathcal{B}}\|\geq \|\vec{\mathcal{A}}\|$.


(3)   By the given conditions,  for all pairs $u\neq v\in [t,\ell]$ with  $u+v=\ell+t$, either $\vec{\mathcal{G}}[\ell]^{(u)}\neq (\emptyset,\ldots, \emptyset)$ and $\vec{\mathcal{G}}[\ell]^{(v)}\neq (\emptyset,\ldots, \emptyset)$, or $\vec{\mathcal{G}}[\ell]^{(u)}=\vec{\mathcal{G}}[\ell]^{(v)}=(\emptyset,\ldots, \emptyset)$.
Since $\vec{\mathcal{G}}[\ell]\neq (\emptyset,\ldots, \emptyset)$, it follows  that either there exist $u\neq v\in [t,\ell]$ with  $u+v=\ell+t$ such that both $\vec{\mathcal{G}}[\ell]^{(u)}$ and $\vec{\mathcal{G}}[\ell]^{(v)}$ are  not equal to  $(\emptyset,\ldots, \emptyset)$, or  $2\mid \ell+t$ and $\vec{\mathcal{G}}[\ell]= \vec{\mathcal{G}}[\ell]^{(\frac{\ell+t}{2})}$. 
If the latter holds, then we are done. We now turn to the former.
From $\vec{\mathcal{G}}[\ell]^{(t)}=(\emptyset,\ldots,\emptyset)$, we obtain $t<u, v<\ell$. 
Combining this with the shiftedness, we obtain that
$\mathcal{G}_k\backslash \mathcal{G}_k[\ell]\neq \emptyset$  for all $k\in[m]$.
For every pair
$u\neq v\in [t,\ell]$ with $u+v=\ell+t$ such that both
$\vec{\mathcal{G}}[\ell]^{(u)}$ and $\vec{\mathcal{G}}[\ell]^{(v)}$ are not equal to
$(\emptyset,\ldots,\emptyset)$,  we choose exactly one of the two boundary layers having
the larger norm; in case of equality, choose either one.
Let $\mathcal C$ be the union of all chosen boundary layers.
Then $\|\vec{\mathcal{G}}[\ell]-\vec{\mathcal{G}}[\ell]^{(\frac{\ell+t}{2})}- \mathcal C\|\leq \|\mathcal C\|$.
Define 
$$
\vec{\mathcal{G}}^*
=
\vec{\mathcal{G}}-\vec{\mathcal{G}}[\ell]+\vec{\mathcal{G}}[\ell]^{(\frac{\ell+t}{2})}+
\{C\setminus\{\ell\}:C\in\mathcal C\},\qquad \vec{\mathcal{B}}=\langle\vec{\mathcal{G}}^*\rangle.
$$

We now verify that $\vec{\mathcal B}$ is shifted. Fix $k\in[m]$
and $1\leq i<j\leq n$. Let $F\in\mathcal B_k$ satisfy
$j\in F$ and $i\notin F$, and put
$
  F'=(F\setminus\{j\})\cup\{i\}.
$
Choose $D\in\mathcal G_k^{*}$ such that $D\subseteq F$. If
$j\notin D$, then $D\subseteq F'$, and hence $F'\in\mathcal B_k$.
We may therefore assume that $j\in D$.
First, suppose that
$
  D\in\mathcal G_k\setminus\mathcal G_k[\ell].
$
Then $\ell\notin D$. Since $\mathcal A_k$ is shifted,
$
  D'=(D\setminus\{j\})\cup\{i\}\in\mathcal A_k.
$
Hence, there exists $E\in\mathcal G_k$ with $E\subseteq D'$. Since
$\ell\notin D'$, we also have $\ell\notin E$, so $E$ is retained in
$\mathcal G_k^{*}$. Thus $E\subseteq F'$ and
$F'\in\mathcal B_k$.
Next, suppose that
$
  D\in\mathcal G_k[\ell]^{(\frac{\ell+t}{2}}.
$
If $j=\ell$, then $D'=(D\setminus\{\ell\})\cup\{i\}$ contains no
$\ell$, and the preceding argument applies. Suppose that $j<\ell$.
Choose $E\in\mathcal G_k$ with $E\subseteq D'$. If $\ell\notin E$,
then $E$ is retained. If $\ell\in E$ and
$|E|=\frac{\ell+t}{2}$, then $E=D'$ and $E$ belongs to the retained middle
boundary layer. Finally, if $\ell\in E$ and
$|E|<\frac{\ell+t}{2}$, choose
$
  h\in D'\setminus E.
$
Then $h<\ell$, and shiftedness gives
$
  (E\setminus\{\ell\})\cup\{h\}\in\mathcal A_k.
$
This set is contained in $D'$ and contains no $\ell$, so it contains
some non-boundary generator belonging to
$\mathcal G_k\setminus\mathcal G_k[\ell]$, which is retained in
$\mathcal G_k^{*}$. Hence again $F'\in\mathcal B_k$.
Finally, suppose that
$
  D=C\setminus\{\ell\}
$
for some generator $C$. Since
$j\in D$ and $i\notin D$, shiftedness of $\mathcal A_k$ gives
$
  C'=(C\setminus\{j\})\cup\{i\}\in\mathcal A_k.
$
Choose $E\in\mathcal G_k$ with $E\subseteq C'$. If $\ell\notin E$,
then $E\subseteq D'$ and $E$ is retained, where
$
  D'=(D\setminus\{j\})\cup\{i\}.
$
If $\ell\in E$ and $|E|=|C|$, then $E=C'$. Since the entire
$|C|$-boundary layer was selected, we have
$
  E\setminus\{\ell\}=D'\in\mathcal G_k^{*}.
$
If $\ell\in E$ and $|E|<|C|$, choose
$
  h\in D'\setminus(E\setminus\{\ell\}).
$
Then
$
  (E\setminus\{\ell\})\cup\{h\}\in\mathcal A_k
$
by shiftedness. This is a subset of $D'$ not containing $\ell$, and
therefore it contains a retained non-boundary generator. Thus
$F'\in\mathcal B_k$ in all cases.

We conclude that each $\mathcal B_k$ is shifted, and hence
$\vec{\mathcal B}$ is a shifted sequence.
By  Lemma \ref{le2}, $\vec{\mathcal{B}}$ is also pairwise cross $t$-intersecting. 
Morover, Lemma \ref{le1} implies
\begin{align*}
\vec{\mathcal{B}}=\|\langle \vec{\mathcal{G}}^*\rangle\| =\|\vec{\mathcal{A}}\|-\|\vec{\mathcal{G}}[\ell]-\vec{\mathcal{G}}[\ell]^{(\frac{\ell+t}{2})}- \mathcal C\|2^{n-\ell}+\|\mathcal C\|2^{n-\ell}\geq \|\vec{\mathcal{A}}\|.
\end{align*}
Thus, we  obtain a non-empty monotone shifted pairwise cross $t$-intersecting sequence $\vec{\mathcal{B}}$  with its extent at most $\ell$ and $\|\vec{\mathcal{B}}\|\geq \|\vec{\mathcal{A}}\|$. 
After deleting redundant members from $\vec{\mathcal G}^{*}$ and keeping the same notation, every $M\in\mathcal G_k[\ell]^{(\frac{\ell+t}{2})}$ remains minimal, since a retained non-boundary generator contained in $M$ would contradict the antichain property of $\mathcal G_k$, while $C\setminus\{\ell\}\subseteq M$ would imply $C\subseteq M$, again contradicting the antichain property of
$\mathcal G_k$. As all other members avoid $\ell$, we obtain $\vec{\mathcal{G}}^*[\ell]= \vec{\mathcal{G}}[\ell]^{(\frac{\ell+t}{2})}$ whenever $\vec{\mathcal B}$ has extent $\ell$.
\end{proof}

\end{document}
